\selectlanguage{english}
\noindent\textbf{Integration, spaces of morphisms, and the choice of weak equivalences.}

These three chapters investigate how much of homotopy theory can be reconstructed from operations on ordinary categories, provided that one keeps the right forms of coherence. Chapter VIII isolates integration for relative $1$-types. Chapter IX returns to the path categories introduced earlier, first to extract an abstract theory of composition and then to compare it with Quillen's constructions. The opening of Chapter X changes the point of departure: instead of beginning with a derivator and its coefficients, Grothendieck asks whether a class of weak equivalences of categories can itself generate the theory. The movement is from a concrete adjunction, through the structure of its spaces of morphisms, to the axioms that make such constructions possible.

The present package includes all of Chapters VIII and IX in the supplied edition, and only the opening of Chapter X. Its thirty-three PDF leaves contain six preliminary leaves and twenty-seven printed pages of mathematical text. Manuscript references must always include the chapter: they run through VIII.1--9, IX.1--25, and X.1--5. Our audit follows those separate sequences. Original statements, including the edition's retained slips and editorial interventions, remain distinguishable from the commentary. The two doubtful transcription readings are resolved by the PDF; neither licenses a mathematical change to Grothendieck's text.

\medskip\noindent\textbf{Relative groupoids and a genuine biadjunction.}

The decisive construction of Chapter VIII is simple enough to write down on an object of the base. For $f:S\to T$ and a category $G$ fibred in groupoids over $S$, set
\[
 (f_!G)_t=\Pi(t\backslash_TG).
\]
The comma category records a point of $G$ together with an arrow from $t$ to its image in $T$; groupoid completion turns this record into a homotopy $1$-type. Its use is not limited to an already fibred category. For any $X\to T$, the same construction gives a relative groupoid $R_TX$, and extension along $X\to R_TX$ is universal for functors into categories fibred in groupoids. Cartesian lifts construct the extension, while their uniqueness supplies the coherence.

This universal property completes the adjunction verification left unfinished in the manuscript. It proves an equivalence of Hom-groupoids, not only a bijection on isomorphism classes. Both triangle comparisons arise from the same extension-and-restriction construction. Transitivity of integration is then the uniqueness of the left biadjoint to composite restriction. The point of the $2$-categorical formulation is thus practical: the automorphisms that would have been discarded by taking components are precisely what make the comparison usable in a further composition. The base-change theorem has the corresponding later expression in the theory of asphericity structures and smooth functors \citep[\S5]{Maltsiniotis2005Asphericite}.

The final descent statements require a more careful separation. The proper or smooth hypotheses of Theorem 3 support its assertions about faithfulness and about the essential image described by coefficients that are \emph{constant on each strict fibre}. They do not justify replacing this condition by the assertion that every vertical transport is an equivalence. A locally constant groupoid family can carry degree-two coherence even when its base is simply connected. We give a finite example over the fourteen-element face poset of the boundary of a tetrahedron. Every fibre is $B(\mathbf Z/2)$ and every transport is the identity, but a nonzero normalized $2$-cocycle prevents the total category from admitting a section. Consequently it cannot be equivalent to a constant family. This disproves part 3 and the subsequent corollary as printed, while leaving the preceding fibrewise-constant criterion intact.

The repair is expressed by higher monodromy: equivalence classes of locally constant groupoid families with a fixed fibre are controlled by maps to its automorphism $2$-type. A $1$-connected base removes ordinary monodromy but can retain the degree-two class; $2$-connectivity removes both \citep[Corollary~2.2.4 and Theorem~2.2.5, arXiv version]{PoleselloWaschkies2005}. A separate, shorter counterexample shows that section groupoids alone do not test all the proposed faithfulness conditions. The correct section tests must also be made after passing to the total category of the first coefficient. Both corrections concern the quantifiers and the coherence in descent, rather than the integration construction itself.

The general faithfulness question also admits an exact answer without smoothness or properness. For each arrow of $T$, form its category of factorizations through objects in the image of $f$. Restriction is faithful on transformations, fully faithful on Hom-groupoids, or an equivalence on Hom-groupoids precisely when these factorization categories are respectively nonempty, connected, or $1$-connected. We prove this three-level criterion by the counit on representables and coherent density, alongside the known untruncated homotopical-epimorphism criterion \citep[Theorem~6.10, arXiv version]{GrothStovicek2016Trees}.

\medskip\noindent\textbf{What composition gives, and what it presupposes.}

Chapter IX begins with a strong idea: composition of paths should explain endpoint functoriality. In a monoidal category $\mathcal H$, Hom objects $C(a,b)$ with associative composition and coherent units determine an ordinary category with morphisms $\operatorname{Hom}_{\mathcal H}(I,C(a,b))$. Inserting these elements into composition produces the bifunctor $C_0^{\mathrm{op}}\times C_0\to\mathcal H$. The proof is formal once the hypotheses are in place. Neither symmetry nor an internal Hom is needed. This is the familiar passage from an enriched category to its underlying ordinary category \citep[\S\S1.1--1.3]{Kelly1982}.

There are two limits to what this argument proves in the passage under review. First, the endpoint maps give isomorphisms only when those endpoint maps are invertible. The early corollary omits that hypothesis, and a discrete two-object example disproves the wording; the manuscript itself restores it on page 8. Second, the abstract argument assumes associative composition on the particular objects $\Ch_\infty(a,b)$. It does not construct that composition. Grothendieck sketches the finite-stage argument for the ind-object $\mathrm{CH}_\infty$ and explicitly leaves the corresponding realized-space verification unfinished.

The possible shortcut through finite tests does not solve this difficulty. The earlier definition of $\Hot_0$ uses homotopy types of finite categories, not just finite CW complexes. The resulting restricted Yoneda functor is not conservative, by the general theorem on small families of unpointed tests, and it is not faithful, by nonconcreteness of the homotopy category \citep[Theorem~2.1]{ArlinChristensen2023} and \citet[pp.~1, 5 of the reprint]{Freyd1970}. Thus equality after passing to those presheaves cannot replace a compatible homotopy on the original mapping spaces. The commentary also distinguishes an ind-object, an ordinary colimit in $\Hot$, and a homotopy colimit. An explicit increasing sequence of finite discrete sets has an ordinary colimit without being essentially constant as an ind-object. A finite-category test such as $BC_2$, on the other hand, is not compact with respect to all filtered homotopy colimits. These examples delimit the arguments that the comparison maps actually require.

Modern localization supplies the intended kind of mapping-space theory. The hammock localization has coherent mapping spaces and a component category equal to the ordinary localization \citep[\S2]{DwyerKan1980Calculating}. This identifies the modern object to which Grothendieck's construction should be compared; it does not by itself prove a comparison for his exact bilateral stabilization and transformation conventions. We keep that model-specific verification explicit. It is not described as a new open problem in the general theory of localization.

\medskip\noindent\textbf{Pointed and linear refinements.}

The later part of Chapter IX asks that the zero map become a specified basepoint in every mapping space. A basepoint in the unpointed homotopy category is insufficient: based homotopies remember data lost under free homotopy. The commentary constructs the pointed morphism associated to a map and a chosen path between basepoints, including its composition law. It also explains the significance of the manuscript's two corrective footnotes. The Cartesian unit in pointed spaces has only one map to every object; the unit that recovers all components is $S^0$ for the smash product. Composition must therefore annihilate either zero input. Forgetting basepoints carries this structure to unpointed composition through a lax monoidal comparison.

Some questions at the level of ordinary localization have particularly direct answers. Initial and terminal objects survive localization without extra hypotheses, by naturality of their unique arrows along both ordinary and inverted generators. Suitable closure under finite products or sums lets the corresponding adjunctions descend. For an additive category, stability of the denominator class under finite direct sums already yields an additive ordinary localization \citep[Lemma~2.2.1]{Krause2021}. The stronger request for complexes of morphisms requires an enhanced setting. We treat complexes of modules explicitly, with connective truncations of derived Hom, derived tensor composition, and the Dold--Kan comparison. This makes the grading visible: positive homotopy groups of the mapping space are negative cohomology groups of the Hom complex.

The suspension and loop objects are then representability statements. Representing the fundamental-group bifunctor yields an adjunction; the equivalences of whole mapping spaces arise from homotopy pushouts and pullbacks. An explicit path complex verifies the loop construction with its differential, its contraction, and its fibre. We also develop a finite counterexample to the proposed force of tests (a) and (b): the full category of pointed sets on one and two elements, with only identities inverted, satisfies both tests and even the displayed mapping-space identities. Its loop and suspension functors are zero, while its ordinary category is nontrivial. The missing strength is therefore not finiteness of the category alone. This calculation is marked as the commentators' development, without a priority claim.

\medskip\noindent\textbf{Localizers and the information retained by coefficients.}

The opening of Chapter X extracts a class $W_{\D}$ by testing maps of indexing categories against constant coefficients. It must be distinguished from the condition that restriction be an equivalence on all coherent diagrams. The localizer axioms express the stability of those constant-coefficient tests under interval homotopy, comma constructions, and integration of fibrations. We complete the reduction from the five displayed axioms to the shorter list, including the role of the twisted-arrow category in proving invariance under opposites. The finite-product and finite-sum questions are then settled by explicit closure arguments and descent of the product and coproduct adjunctions. No model structure is needed for these formal consequences.

Realization is a further theorem. For an accessible fundamental localizer, the later model structures and derivator constructions give a coefficient theory that realizes its weak equivalences \citep[Theorems~4.2.15 and 5.2.15]{Cisinski2006}; \citep[Theorem~6.11 and its dual]{Cisinski2003}. This does not reconstruct an arbitrary original derivator from $W_{\D}$. The represented derivators of sets and of abelian groups have the same component-detecting localizer, although their coefficient categories are inequivalent. The represented truth-value lattice gives an equally concrete nonzero example whose localizer detects only emptiness, disproving any attempt to impose component detection on all nonzero derivators. These examples show precisely what Chapter X has isolated: the homotopy shapes visible to a theory, while its further algebra of coefficients remains additional structure.

\nocite{*}
